% Supplement: results this repository holds that the submitted paper does not.
%
% Every number here is a macro from frozen_numbers.tex, so this file is
% regenerated with the rest of the manuscript and cannot drift from the
% artifacts on its own.
%
% Why this file exists. The PSB submission, "Auditable Algebraic Counting Field
% for Cryptic-Pocket Detection from Apo Structures", names commit 5d1726a as its
% frozen snapshot and directs the reader to audit the chronology from the commit
% graph of this repository. A reader who follows that pointer arrives at results
% the paper does not mention, including an architecture that is not separable
% from pLM-NN where the paper reports a resolved deficit. Nothing here
% contradicts the paper -- the paper's claim is about the counting field and
% remains true of it -- but without a statement of which document is
% authoritative, a reader cannot tell whether they have found a stronger result
% or an inconsistency. This file says which.
%
% tools/check_submission_consistency.py enumerates the divergence, and its
% artifact is results/official_fold/SUBMISSION_CONSISTENCY.json. That tool
% checks 47 of the submission's printed numbers against the macros that
% generated them; all 47 agree, which is why this supplement is about process
% counts and scope rather than about results.

\section*{Supplement: results beyond the submitted paper}
\label{sec:supplement}

The submitted paper is authoritative for every claim it makes. This supplement
holds work in the repository that postdates or exceeds it, and exists so that a
reader arriving from the paper is not left to guess which document governs.

Every bracketed pair below is a 95\% confidence interval from a paired
nonparametric bootstrap over units, the same procedure and the same resample
count the paper uses, and every difference is computed on the intersection where
both methods are defined. A difference whose interval contains zero is called a
point-estimate lead or deficit and never a win or a loss, whatever its sign.

\paragraph{Two process counts in the paper are wrong, and both understate us.}
The submission's Limitations state ``12 indexed test reads and 12 scored
architectures''. Neither figure is what this repository reports.

The ledger recorded \textbf{13} indexed reads at the submission commit itself,
so that number was already off by one when the paper was frozen; it is 13 today.
The architecture count was correct at \texttt{5d1726a} and is now
\NArchOnFold{}. The jump is not new work alone. The ledger had three signals for
``this artifact took a number off the held-out fold'', and all three keyed on
how an artifact \emph{spelled} its metric, so eighteen artifacts that store each
method's score under the method's own name were invisible to it. The fix added a
fourth signal: the \texttt{reads\_test\_fold} declaration that every artifact in
this repository is already required to carry, which is the only one of the four
independent of vocabulary.

Both corrections widen the multiplicity that the paper's small official-fold
margins should be read against, so both cut against the paper's own numbers
rather than for them. The full accounting is \NFoldAccesses{} accesses to the
fold, of which \NIndexedReads{} carry an explicit read index and
\NLineageReads{} name the lineage of the detector the paper deploys.

\paragraph{Six architectures of ours outscore the one the paper evaluates.}
The paper evaluates a single detector of ours, the counting field, and reports
it as trailing pLM-NN by \PlmAucDelta{} \PlmAucCI{} on the official fold. That
remains true of that detector, and the paper's Discussion is correct to bound
the accuracy claim by it. On the same units, in a single pass over one residue
universe, \GrandNAboveTab{} of our \GrandNOurs{} architectures have a higher
mean per-unit ROC-AUC than it does.

That last sentence is an ordering of point estimates and nothing more. No
paired interval was computed between our own architectures, they share a feature
set and are strongly correlated, and \NArchOnFold{} of them have now read this
fold, so the ordering among them carries no information about which is better;
it is reported only to establish that the deployed detector is not the highest
of them. The comparison that does carry an interval is against the baseline:
the best of ours reaches \GrandBestRoc{} against pLM-NN's \GrandPlmRoc{}, a
paired difference of \GrandDPlmRoc{} \GrandDPlmRocCI{} on
\GrandDPlmRocNPaired{} units, ahead on \GrandDPlmRocWin{} and behind on
\GrandDPlmRocLoss{} --- an interval containing zero, and therefore a
point-estimate lead and not a win. On PR-AUC the point estimate changes sign,
\GrandDPlmPr{} \GrandDPlmPrCI{}, and that interval contains zero too. Both
metrics say parity; neither says we are ahead.

This is a development read on a fold already worked too hard to settle
anything. It is not offered as a correction to the paper's claim; it is offered
because a reader of the repository will find it, and should find it labelled.

The section that follows was written for the long-form manuscript and is
reproduced here unchanged.

\subsection*{All three baselines in one pass, and the metric that disagrees}
\label{sec:grand}
The comparisons in the main text are pairwise and were taken at different times,
which leaves two questions a reader cannot answer from them. Was every method
scored on the same residues? And does a method that leads on ROC-AUC also lead
on the metric that weights a sparse positive class? This section closes both by
scoring \GrandNMethods{} methods in one pass over one residue universe on all
\GrandN{} units: \GrandNOurs{} of ours --- the deployed counting field, the
geometry and nonlocal-seam families at both gate radii, and their two equal-$z$
fusions --- together with P2Rank 2.5.1, CryptoBench's own pLM-NN readout and
PocketMiner. Every paired row below has
$n_{\text{paired}} = \GrandDPlmRocNPaired{}$, so no difference here
mixes \emph{which structures were scorable} into a number read as \emph{which
method is better}. The reason to do this in a single pass rather than to
assemble it from the reads already taken is that a cell missing now cannot be
filled in later: the universe would have to be rebuilt, and the earlier numbers
would no longer be commensurable with the new ones.

The pass is complete in its columns and not in its rows, and the distinction
matters. Every method named above is scored on every unit, so no cell of the
table is missing. But \GrandNOurs{} is not all \NArchOnFold{} of the
architectures that have taken a number off this fold; the others --- quotient
banks, the algebraic fields, \texttt{threshold\_field}, the remaining seam
variants --- were read in earlier probes and are not re-scored here. That is a
deliberate refusal rather than an omission. Re-running them would be another
sweep of the held-out fold, and a section whose argument is that this fold has
already been worked too hard to settle anything cannot coherently work it
harder to tidy its own table. Their readings stand where they were taken, in
the ledger, with the multiplicity they imply already charged against every
margin below.

Three terms are new. The \textbf{nonlocal seam} family measures, per residue,
whether a position is packed in the core while carrying long-range contact
support and conformational budget in its neighbourhood on a peelable depth
layer --- the configuration of a site that can open without presenting as a
surface pocket in the apo structure. It was written because
\texttt{geometry\_field} reaches fold-mean parity while anti-ranking buried
cryptic residues on a handful of chains, and it counts what the surface-facing
families do not. An \textbf{equal-$z$ fusion} of two detectors replaces each
score by its within-chain standardised value and adds the two, unweighted; no
coefficient is fitted, and ``equal'' is the entire rule, so the fusion inherits
the property that nothing in the inference path was tuned on any fold. The
\textbf{$r{=}14$ gate} is the spatial gate applied at 14~\AA{} rather than the
deployed 18~\AA{}. That radius was not chosen here: it was selected on twelve
cluster-disjoint halves of the \emph{training} fold, with each half of the
splits independently choosing the same arm, at a paired $+0.0025$ uniform across
all four pocket-size strata. Scoring it is one read of a decision already made,
and no other radius is scored on this fold.

The best of our architectures under this read is \GrandBest{}. On mean per-unit
ROC-AUC it reaches \GrandBestRoc{} against \GrandPlmRoc{} for pLM-NN, a paired
difference of \GrandDPlmRoc{} with interval \GrandDPlmRocCI{}, ahead on
\GrandDPlmRocWin{} units and behind on \GrandDPlmRocLoss{}. \textbf{On mean
per-unit PR-AUC the ordering reverses}: \GrandBestPr{} against \GrandPlmPr{}, a
paired difference of \GrandDPlmPr{} with interval \GrandDPlmPrCI{}, ahead on
\GrandDPlmPrWin{} and behind on \GrandDPlmPrLoss{}. Both intervals contain zero.
The honest statement is parity in both directions, and the two metrics
disagreeing about the sign is the more informative half of it: PR-AUC weights
the sparse positive class, so pLM-NN's advantage is concentrated where the
labelled pocket is small, which is consistent with the length stratification in
the main text and with a sequence model carrying context that local geometry
does not. Reporting only the ROC-AUC row would have converted a null into a
lead.

Against the other two baselines the lead is resolved on both metrics.
\GrandBest{} exceeds P2Rank by \GrandDPTwoRoc{} \GrandDPTwoRocCI{} on ROC-AUC
and \GrandDPTwoPr{} \GrandDPTwoPrCI{} on PR-AUC, and PocketMiner by
\GrandDPmRoc{} \GrandDPmRocCI{} and \GrandDPmPr{} \GrandDPmPrCI{}. For
comparison on the same units, the deployed \texttt{table\_field} sits at
\GrandTabRoc{} and \GrandTabPr{}; the geometry and seam families are what move
the standing against pLM-NN from a resolved deficit to a null, and they do not
move it past one.

None of this is confirmatory. \NArchOnFold{} of our architectures have now taken
a number off this fold, and a margin of a hundredth on \GrandN{} units is well
inside the multiplicity that implies; the fold's standard error on a single
mean is $0.0161$. The confirmatory accuracy claim of the paper remains the one
made for \texttt{table\_field} against P2Rank on the external set, and the
spent external sets are not re-opened to manufacture a second one for these
architectures. What the read establishes is narrower and still worth stating:
on one residue universe, with no cell missing, a detector carrying no learned
encoder and no floating-point discriminant is not separable from a supervised
network over a 3-billion-parameter protein language model, and is separably
ahead of both published pocket detectors. Figures~\ref{fig:grand-standing} and
\ref{fig:grand-paired} show the standing and the paired differences; in the
latter, a bar whose interval crosses zero is drawn grey whatever the sign of its
point estimate, so that a parity cannot be read as a win.

\begin{figure}[t]
\centering
\includegraphics[width=\textwidth]{fig_grand_baseline_standing.png}
\caption{\FigCapGrandBaselineStanding{}}
\label{fig:grand-standing}
\end{figure}

\begin{figure}[t]
\centering
\includegraphics[width=\textwidth]{fig_grand_baseline_paired.png}
\caption{\FigCapGrandBaselinePaired{}}
\label{fig:grand-paired}
\end{figure}
